\documentclass[10pt,a4paper]{amsart}
\usepackage[margin=19mm]{geometry}
\usepackage{amsmath,amssymb,mathtools}
\usepackage[scr=rsfs,cal=euler]{mathalfa}
\usepackage{xfrac}
\usepackage[unicode,hidelinks,bookmarks=false]{hyperref}
\newcommand{\ie}{i.e.\ }
\newcommand{\cf}{cf.\ }
\newcommand{\resp}{respectively\ }
\newcommand{\Subsection}{\subsection}
\providecommand{\nobreakdash}{\nobreak\hbox{-}}
\setlength{\emergencystretch}{2em}
\multlinegap0pt
\usepackage{bm}
\usepackage{bbm}
\usepackage{dsfont}
\usepackage{xspace}
\usepackage{cancel}
\usepackage{mathtools}
\usepackage{amsthm}
\makeatletter
\@ifundefined{thm}{\newtheorem{thm}{Theorem}}{}
\@ifundefined{conj}{\newtheorem{conj}[thm]{Conjecture}}{}
\@ifundefined{pr}{\newtheorem{pr}[thm]{Proposition}}{}
\makeatother
\newcommand{\E}{\mathbb{E}}
\title[Second-Order Schalkwijk-Kailath Coding for Autoregressive Ga]{Second-Order Schalkwijk-Kailath Coding for Autoregressive Gaussian Channels}
\author{Jun Su, Guangyue Han, Shlomo Shamai (Shitz)}
\date{Source: arXiv:2601.09329v3 (CC BY 4.0)}
\begin{document}
\maketitle

\noindent\emph{Source and licence.} Second-Order Schalkwijk-Kailath Coding for Autoregressive Gaussian Channels. Version arXiv:2601.09329v3 (CC BY 4.0), distributed under the Creative Commons Attribution 4.0 International licence, as recorded in the arXiv licence metadata for this exact version and shown on the abstract page https://arxiv.org/abs/2601.09329v3. Full rights notice: licenses/arxiv-2601.09329-CC-BY-4.0.txt.\par\smallskip

\noindent\textit{Editorial scope.} Counted unit: the Pr whose source label is \texttt{prop:interleaved-AR1} in the article (source0.tex lines 1212--1223, source label \texttt{prop:interleaved-AR1}), delivered together with the article's own complete original proof of it (source0.tex lines 1225--1368, one \texttt{proof} environment of 144 lines). No mathematics is written, repaired or re-derived: statement and proof are the article's own text line for line. Required context retained uncounted: channel-arp (lines 106--108); ISK1-thm (lines 235--248); Butman-conj:label (lines 439--444); SK2-coding-input (lines 460--462); SK2-coding-message (lines 464--472); main-theorem-sk2 (lines 550--571); power-eqn-sk2-thm (lines 560--562). Every definition, display and label of those lines is retained unchanged. Omitted from this package and not redistributed: 1--105, 109--234, 249--438, 445--459, 463--463, 473--549, 563--1211, 1224--1224, 1369--1625. Nothing inside a retained excerpt is omitted or altered: every retained body is delivered exactly as the article's lines are written, so no omission receipt of any kind accompanies this package. Citations: the retained excerpts carry no citation command at all, so no bibliography entry of the article is transcribed and no reference is redistributed. Reference adaptation: none was needed and none was made; every reference inside the delivered unit resolves inside it, so no unresolved reference or citation is produced. Printed slips kept literal: no printed word, symbol, hypothesis or number of the counted statement or of its proof was altered. Layout and class: the article's own class is \texttt{article} with packages including babel, geometry, amsmath, graphicx, yfonts, amssymb, amsthm, bm, comment, bbm, mathtools, hyperref; the wrapper is the lane's pinned \texttt{amsart} editorial wrapper at 10pt on A4, which restates the theorem-like environments the retained text needs and adds the article's own macro definitions that the retained text needs (\texttt{\textbackslash E}), with extra wrapper packages (bm, bbm, dsfont, xspace, cancel, mathtools). The delivered numbering is the unit's own. Assets: no retained span contains a figure environment, a graphics-inclusion command, a PDF-page inclusion, a table or a TikZ picture; no image file is copied into this package, the package declares no asset dependency and no image file is redistributed with it. Licence: version arXiv:2601.09329v3 of this article is distributed under the Creative Commons Attribution 4.0 International licence, as recorded in the arXiv licence metadata for this exact version and shown on the abstract page https://arxiv.org/abs/2601.09329v3. A full rights notice is delivered at licenses/arxiv-2601.09329-CC-BY-4.0.txt. Characters: every retained excerpt is pure ASCII, so the delivered engine, XeLaTeX with its default Latin Modern text font, reports no missing character.
\begin{equation}\label{channel-arp}
Y_i = X_i + Z_i, \quad i\ge 1,
\end{equation}

\begin{thm}\label{ISK1-thm}
Suppose the noise $\{ Z_i \}$ is an $AR(p)$ Gaussian process with parameters $\beta_k$, $\vert \beta_k\vert<1$ for all $k=1,2,...,p$, namely, it has the power spectral density
\begin{equation*}
S_Z(e^{i\theta}) = \vert H_Z(e^{i\theta}) \vert^2 = \frac{1}{\vert L_Z(e^{i\theta}) \vert^2} = \frac{1}{\vert\prod_{k=1}^{p}(1+\beta_ke^{i\theta})\vert^2}.
\end{equation*}
Then, we have
\begin{equation}\label{I-SK1-thm}
C_\mathrm{SK1}(P)= \max_{\gamma} \log  |\gamma|,
\end{equation}
where the maximum is taken over all $\gamma\in\mathbb{R}$ with $|\gamma|>1$ satisfying the power constraint 
\begin{equation}\label{power-eqn-sk1-thm}
   (\gamma^2-1)H^2_Z(\gamma^{-1}) = P.
\end{equation}
\end{thm}

\begin{conj}[Corrected form of Butman's conjecture]\label{Butman-conj:label}
For stationary AR($p$) Gaussian channels \eqref{channel-arp},
\[
    C_{\mathrm{FB}}(P)=C_{\mathrm{SK1}}(P).
\]
\end{conj}

\begin{equation}\label{SK2-coding-input}
X_i = V_i - \E{[V_i\mid Y_1^{i-1}]}, \quad i=3,4,\dots
\end{equation}

\begin{equation}\label{SK2-coding-message}
\left\{ 
\begin{aligned}
V_{i+1} &= ~aV_i + bV_{i-1},\quad i\ge2,\\
V_2 ~~&= ~U_2,\\
V_1 ~~&= ~U_1,
\end{aligned}
\right.
\end{equation}

\begin{thm}\label{main-theorem-sk2}
Suppose the noise $\{ Z_i \}$ is an $AR(p)$ Gaussian process with parameters $\beta_k$, $\vert \beta_k\vert<1$ for all $k=1,2,...,p$, namely, it has the power spectral density
\begin{equation*}
S_Z(e^{i\theta}) = \vert H_Z(e^{i\theta}) \vert^2 = \frac{1}{\vert L_Z(e^{i\theta}) \vert^2} = \frac{1}{\vert\prod_{k=1}^{p}(1+\beta_ke^{i\theta})\vert^2}.
\end{equation*}
Then, we have
\begin{equation}\label{I-SK2-thm}
C^\mathrm{g}_{\mathrm{SK2}}(P) =\sup_{\gamma_1,\gamma_2} \log  |\gamma_1\gamma_2|
\end{equation}
and the supremum is taken over all $\gamma_1,\gamma_2\in\mathbb{C}$ satisfying $|\gamma_1|\wedge|\gamma_2|>1$, $\gamma_1+\gamma_2\in \mathbb{R}$, $\gamma_1\gamma_2\in\mathbb{R}$, and the power constraint 
\begin{equation}\label{power-eqn-sk2-thm}
\frac{1}{\Delta}\left( \frac{L_{11}}{\gamma^2_2}+\frac{L_{22}}{\gamma_1^2} +\frac{2L_{12}}{\gamma_1\gamma_2}  \right)= P,
\end{equation}
where 
\begin{equation}\label{main-thm-notations}
\begin{aligned}
\Delta &= (\gamma_1-\gamma_2)^2(L_{11}L_{22}-L_{12}^2),\\
L_{ij} &=\frac{c_{i}c_{j}L_Z(\gamma_i^{-1})L_Z(\gamma_j^{-1})}{\gamma_i\gamma_j-1} \quad \mbox{for $i,j=1,2,$}\\
c_1 &= \frac{1}{\gamma_1(\gamma_1-\gamma_2)}, \quad c_2 = \frac{1}{\gamma_2(\gamma_2-\gamma_1)}.
\end{aligned}
\end{equation}
\end{thm}

\begin{pr}
\label{prop:interleaved-AR1}
Consider the stationary AR(2) Gaussian channel
$$
Y_i=X_i+Z_i,\qquad Z_i+\beta Z_{i-2}=W_i,\quad i\in\mathbb Z,
$$
where \(0<|\beta|<1\) and \(\{W_i\}\) is a zero-mean, unit-variance white Gaussian process. Then, for every $P>0$,
$$
C_{\mathrm{SK2}}(P) = C_{\mathrm{FB}}(P) > C_{\mathrm{SK1}}(P) = R_{\mathrm{Butman}}(P).
$$
Consequently, Conjecture~\ref{Butman-conj:label} is false for stationary AR(2) Gaussian channels.
\end{pr}

\begin{proof}
Define $Z_k^{(e)}=Z_{2k}$ and $Z^{(o)}_k=Z_{2k-1}$ for $k\in\mathbb{Z}$. Since the noise process $\{Z_i\}$ satisfies $Z_i+\beta Z_{i-2}=W_i$, we have
$$
Z^{(e)}_{k}+\beta Z_{k-1}^{(e)}=W_{2k}, \qquad Z^{(o)}_{k}+\beta Z^{(o)}_{k-1}=W_{2k-1}.
$$
Therefore, $\{Z^{(e)}_k\}$ and $\{Z^{(o)}_k\}$ are two independent stationary AR(1) Gaussian processes with parameter $\beta$, driven by the independent white Gaussian noises $\{W_{2k}\}$ and $\{W_{2k-1}\}$, respectively. Consequently, the original AR(2) Gaussian channel can be viewed as two independent parallel AR(1) Gaussian channels obtained by separating the even and odd channel uses.

Let $C_{\mathrm{FB}}(P;\beta)$ denote the feedback capacity of either AR(1) subchannel. Let $P_e$ and $P_o$ denote the average input powers per use of the even and odd AR(1) subchannels, respectively. Since each subchannel occupies one half of the original channel uses, the overall average-power constraint becomes
$$
\frac{P_e+P_o}{2}\le P.
$$
Consequently, the feedback capacity of the original channel is
\begin{equation*}\label{CFB-interleave-ar2}
\begin{aligned}
C_{\mathrm{FB}}(P)
&= \max_{\frac{P_{\mathrm e}+P_{\mathrm o}}{2}\le P}
\left\{
\frac12 C_{\mathrm{FB}}(P_{\mathrm e};\beta) + \frac12 C_{\mathrm{FB}}(P_{\mathrm o};\beta) \right\}.
\end{aligned}
\end{equation*}
Since $C_\mathrm{FB}(P;\beta)$ is concave and non-decreasing in $P$, we have
$$
\frac12 C_\mathrm{FB}(P_{\mathrm e};\beta) + \frac12 C_{\mathrm{FB}}(P_{\mathrm o};\beta)
\le
C_{\mathrm{FB}}\left(\frac{P_{\mathrm e}+P_{\mathrm o}}{2};\beta\right)
\le
C_{\mathrm{FB}}(P;\beta),
$$
with equality attained by $P_{\mathrm e}=P_{\mathrm o}=P$. Hence
\begin{equation}\label{CFB-interleave-ar2-2}
C_{\mathrm{FB}}(P)= C_{\mathrm{FB}}(P;\beta)=\log x_0,
\end{equation}
where $x_0>1$ is the unique solution of
\begin{equation}\label{eq:interleaved-capacity}
P= \frac{x_0^2-1} {\left(1+|\beta|x_0^{-1}\right)^2}.
\end{equation}

Now, applying the SK(2) coding scheme \eqref{SK2-coding-input}--\eqref{SK2-coding-message} and choosing the characteristic pair $(\gamma_1,\gamma_2)$ as
\begin{equation}\label{chara-pair-interleave-ar2}
(\gamma_1,\gamma_2) =
\begin{cases}
\left(\sqrt{x_0},-\sqrt{x_0}\right), & \beta>0,\\[1mm]
\left(\mathrm{i}\sqrt{x_0},-\mathrm{i}\sqrt{x_0}\right), & \beta<0.
\end{cases}
\end{equation}
Then
$$
\gamma_1+\gamma_2=0, \qquad \gamma_1\gamma_2=-\operatorname{sign}(\beta)x_0.
$$
In particular, 
$$
|\gamma_1|=|\gamma_2| = \sqrt{x_0}>1,
$$
and both $\gamma_1+\gamma_2$ and $\gamma_1\gamma_2$ are real. Moreover, 
$$
b = -\gamma_1\gamma_2=\operatorname{sign}(\beta)x_0\neq 0.
$$
Thus, it remains to verify that the characteristic pair $(\gamma_1,\gamma_2)$ satisfies the power constraint in Theorem~\ref{main-theorem-sk2}.
%and the corresponding message recursion is
% $$
% V_{i+1}=\operatorname{sign}(\beta)x_0V_{i-1}.
% $$
% Thus, defining $V_k^{(o)}=V_{2k-1}$ and $V_k^{(e)}=V_{2k}$ for $k\ge 1$, we have
% $$
% \begin{cases}
% V_{k+1}^{(o)} = \operatorname{sign}(\beta)x_0V_{k}^{(o)}\\[4pt]
% V_{1}^{(o)} = U_1
% \end{cases},
% \qquad
% \begin{cases}
% V_{k+1}^{(e)} = \operatorname{sign}(\beta)x_0V_{k}^{(e)}\\[4pt]
% V_{1}^{(e)} = U_2
% \end{cases}.
% $$
%We next verify that this the characteristic pair $(\gamma_1,\gamma_2)$ satisfies the power constraint in Theorem~\ref{main-theorem-sk2}.
For both $\beta>0$ and $\beta<0$, 
$$
\gamma_1^2=\gamma_2^2=\operatorname{sign}(\beta) x_0,
$$
and, since $L_Z(z)=1+\beta z^2$,
$$
L_Z(\gamma_1^{-1})=L_Z(\gamma_2^{-1}) = 1 +\frac{\beta}{\operatorname{sign}(\beta) x_0}=1 +\frac{|\beta|}{ x_0}.
$$
Moreover,
$$
c_1=c_2=\frac{1}{2\operatorname{sign}(\beta) x_0}.
$$
Substituting these expressions into the power constraint \eqref{power-eqn-sk2-thm} gives
$$
\frac{x_0^2-1}{(1 +|\beta| x^{-1}_0)^2}=P,
$$
where the equality follows from \eqref{eq:interleaved-capacity}. Thus, the chosen characteristic pair $(\gamma_1,\gamma_2)$ in \eqref{chara-pair-interleave-ar2} is feasible in Theorem~\ref{main-theorem-sk2}. Consequently,
$$
C^\mathrm{g}_\mathrm{SK2}(P)\ge \log|\gamma_1\gamma_2|=\log x_0.
$$
Since $C_\mathrm{SK2}(P)\ge C^\mathrm{g}_\mathrm{SK2}(P)$, together with \eqref{CFB-interleave-ar2-2}, it holds that
$$
 C_{\mathrm{FB}}(P)=\log x_0 \le C_{\mathrm{SK2}}(P) \le C_{\mathrm{FB}}(P),
$$
and hence
\begin{equation}\label{CSK2-CFB-AR2-Inter}
C_\mathrm{SK2}(P)=C_\mathrm{FB}(P) = \log x_0.
\end{equation}

It remains to prove that $C_{\mathrm{SK1}}(P)<C_{\mathrm{FB}}(P)$. To see this, Theorem~\ref{ISK1-thm} gives
\begin{equation}\label{CFB-sk1-inter}
C_\mathrm{SK1}(P)=\max\log |\gamma|,
\end{equation}
where the maximum is taken over all $\gamma\in\mathbb{R}$ with $|\gamma|>1$ satisfying 
\begin{equation}\label{eq:SK1-interleaved}
P= \frac{\gamma^2-1} {\left(1+\beta \gamma^{-2}\right)^2}.
\end{equation}
Let $\gamma_0$ be an optimizer of \eqref{CFB-sk1-inter} and define
$$
\Phi_\beta(x) = \frac{x^2-1}{\left(1+|\beta| x^{-1}\right)^2}, \qquad x>1.
$$
Since
$$
\Phi_\beta'(x) = \frac{2x\left(x^3+2|\beta| x^2-|\beta|\right)} {(x+|\beta|)^3} >0, \qquad x>1,
$$
the function $\Phi_\beta$ is strictly increasing. Moreover, 
$$
1+\frac{|\beta|}{|\gamma_0|} > 1+\frac{\beta}{\gamma_0^2} >0.
$$
Using \eqref{eq:SK1-interleaved}, we obtain
$$
\Phi_\beta(|\gamma_0|) < \frac{\gamma_0^2-1} {\left(1+\beta \gamma_0^{-2}\right)^2} = P.
$$
On the other hand, \eqref{eq:interleaved-capacity} gives
$$
\Phi_\beta(x_0)=P.
$$
The strict monotonicity of $\Phi_\beta$ implies $|\gamma_0|<x_0$. Consequently
$$
C_{\mathrm{SK1}}(P) = \log |\gamma_0| < \log x_0 = C_{\mathrm{FB}}(P),
$$
as desired.

Finally, it follows from $R_{\mathrm{Butman}}(P) = C_{\mathrm{SK1}}(P)$ and \eqref{CSK2-CFB-AR2-Inter} that 
$$
R_{\mathrm{Butman}}(P) < C_{\mathrm{SK2}}(P) = C_{\mathrm{FB}}(P),
$$
which disproves Conjecture~\ref{Butman-conj:label}. This completes the proof. 
\end{proof}

\end{document}
